% !TeX root = ../paper.tex
\section{Ablation Protocol Details}
\label{appendix:ablation-details}

The standard component ablations use SmartVLA-Pi and the same
target-domain source trajectories, adaptation update count,
random seeds, evaluation instances, and trajectory-level success
contract as the full SmartVLA-Pi model. Table~\ref{tab:ablation-settings}
defines the settings evaluated in Table~\ref{tab:ablation}.

The ablations address three aspects of the training framework.
Stage-removal experiments evaluate the contribution of each
stage to the complete training procedure, including the training
exposure associated with that stage. Data-construction controls
examine the choices used to form local supervision: uniform
segmentation changes the temporal boundaries, while removing
grade conditioning changes the applicability rules for instruction
families. The raw-data adaptation control tests whether local
instruction supervision remains important during target-domain
adaptation. The fine-grained semantic-only setting examines whether
enriching whole-trajectory language is sufficient to obtain the
instruction-following gains of the full framework. It applies
fine-grained rewrites to complete trajectories during pretraining
and adaptation, without warmup, action-aligned local segmentation,
or visual-prompt supervision. This setting implements a
FineVLA-inspired supervision strategy within our experimental
pipeline.

% !TeX root = ../paper.tex
\begingroup
\normalsize
\setlength{\tabcolsep}{4pt}
\setlength{\LTcapwidth}{\linewidth}
\setlength{\LTpre}{\baselineskip}
\setlength{\LTpost}{\baselineskip}
\renewcommand{\arraystretch}{1.10}
\begin{longtable}{@{}
  >{\raggedright\arraybackslash}p{(\linewidth - 2\tabcolsep) * \real{0.25}}
  >{\raggedright\arraybackslash}p{(\linewidth - 2\tabcolsep) * \real{0.75}}@{}}
\caption{Definitions of the SmartVLA-Pi ablation settings.}\label{tab:ablation-settings}\\
\toprule
Setting & Removed or changed component \\
\midrule
\endfirsthead
\toprule
Setting & Removed or changed component \\
\midrule
\endhead
\bottomrule
\endfoot
\bottomrule
\endlastfoot
Full SmartVLA-Pi & warmup, pretraining, and adaptation with all steering-oriented supervision \\
\midrule
w/o warmup & Skip image-free warmup \\
\midrule
w/o pretraining & Skip pretraining and adapt from the warmup checkpoint \\
\midrule
From raw VLM & Skip SmartVLA's warmup and pretraining and begin target-domain training from the original VLM weights; retain the full method's adaptation data processing and steering supervision \\
\midrule
Raw-data adaptation & Retain the warmup and pretraining weights, but fine-tune directly on the corresponding benchmark's original demonstrations and task instructions, without action-aligned segmentation, local-instruction rewriting, or spatial visual annotation \\
\midrule
Uniform segmentation & Replace action-aligned segmentation with fixed-interval temporal cuts, approximately matching the full method's segment count and mean length; retain the subsequent annotation and training procedures \\
\midrule
w/o grade conditioning & Retain the three rewrite families but remove the High/Mid/Low rules governing which families are used; apply a common rewriting rule while retaining semantic and visual consistency review \\
\midrule
Fine-grained sem. only & Two-stage pretraining and adaptation using fine-grained language rewrites on complete trajectories in both stages; no warmup, action-aligned local segmentation, or visual-prompt supervision \\
\end{longtable}
\endgroup
%

The two data-construction controls each change only the specified
component and retain original complete-task supervision and the
original-to-local supervision sampling ratio. Uniform temporal
segmentation approximately matches the full method's segment
count and mean length, with subsequent annotation and training
procedures retained. Its average success rate of 7.45\%, compared
with 14.67\% for the full model, supports the importance of
aligning segment boundaries with control phases. Removing grade
conditioning yields 12.25\%, supporting the use of
interaction-dependent instruction families.

The fine-grained semantic-only setting reaches 7.00\% on average,
showing that whole-trajectory language enrichment alone does not
recover the full framework's gains in this setting. The temporal
segmentation and grade-conditioning controls provide more
targeted evidence for the roles of local alignment and rewriting
applicability. Together, these comparisons support the design
of action-aligned, grade-conditioned supervision for local
instruction following.

